% problem-000451 3 DulongPetit原子量估算
\section*{3. DulongPetit原子量估算}
\textit{下列为分问。}

\subsection*{3-1}
早在19世纪初期，法国科学家Dulong和Petit测定⽐热时，发现⾦属的⽐热 $\left( c _ { m } \right)$ 与其原⼦量的乘积近似为常数 $6 \operatorname { c a l } \mathrm { g } ^ { - 1 } \ { } ^ { \circ } \mathrm { C } ^ { - 1 } \left( 1 \operatorname { c a l } = 4 . 1 8 \operatorname { J } \right)$ )。当时已知的原⼦量数据很少，因此，可利⽤⽐热推算原⼦量，进⽽采⽤其他⽅法分析得到更精确的原⼦量。

\subsection*{3-1}
-1 将 $4 0 . 0 \mathrm { g }$ ⾦属M块状样品加热到 $1 0 0 \ ^ { \circ } \mathrm { C } ,$ ,投⼊ 50.0 g 温度为 $1 5 . 2 ~ ^ { \circ } \mathrm { C }$ 的⽔中，体系的温度为 $1 7 . 2 ~ ^ { \circ } \mathrm { C } _ { \circ }$ 推算该⾦属的摩尔质量。

\subsection*{3-1}
-2 取⾦属M的粉末样品1.000 g，加热与氧⽓充分反应，得氧化物1.336 ${ \mathrm { g } } _ { \circ }$ 。计算该⾦属的摩尔质量，推测其氧化物的化学式（⾦属与氧的⽐例为简单整数⽐）。

\subsection*{3-1}
-3M是哪种⾦属？

\subsection*{3-2}
电解法⽣产铝须⽤纯净的氧化铝。铝矿中常含⽯英、硅酸盐等杂质，需预先除去。在拜耳法处理过程中，硅常以硅铝酸盐 $( \mathrm { N a _ { 6 } A l _ { 6 } S i _ { 5 } O _ { 2 2 } { \cdot } 5 H _ { 2 } O } )$ “泥”的形式沉积下来。现有⼀种含10.0 wt%⾼岭⼟ $\mathrm { . ( A l _ { 2 } S i _ { 2 } O _ { 7 } { \cdot } 2 H _ { 2 } O ) }$ 的⽔铝⽯ $\mathrm { \cdot { A l } ( O H ) _ { 3 } ] }$ 原料，计算纯化处理中铝的损失率。
